すべての累積世代の高被引用論文に対しこの判定を行った結果、536論文中77論文がオブソレッセンス状態にあると判定された。オブソレッセンス状態にあると判定された、および判定されなかった高被引用論文の、各指標値を図\ref{fig:arr1}に示す。オブソレッセンス状態にあると判定された論文では出版からピークまでの期間のモードは4年、ピークから被引用の十分な減衰までの期間のモードは7年、出版から被引用の十分な減衰までの期間のモードは12年であった。この結果は、P.D.B Paroloらの報告（生物分野では出版後5年以内にピークが訪れる）を支持している。一方、オブソレッセンス状態にない論文では出版からピークまでの期間は12年と大きく異なっていた。オブソレッセンス状態にあると判定された高被引用論文における被引用数のプロットを図\ref{fig:gr7}に示す。出版から被引用が十分に減衰するまでの期間は最小で10年、最大で84年であった。同図からは近年の論文ほど出版直後に高い被引用を受ける傾向があると推察される。
